\documentclass[a4paper,oneside,12pt]{amsart}

%\usepackage[spanish, activeacute]{babel}
\usepackage[utf8]{inputenc}
%\usepackage[latin1]{inputenc}
\usepackage{latexsym}
\usepackage{multicol}
\usepackage{enumerate}
\usepackage{enumitem}
\usepackage[heightrounded]{geometry}
\geometry{top=2cm,bottom=2cm,left=2cm,right=2cm}
\usepackage{graphicx}

\DeclareMathOperator{\cond}{cond}
\DeclareMathOperator{\Id}{Id}
\DeclareMathOperator{\re}{Re}
%\DeclareMathOperator{\im}{Im}
\newcommand{\I}{\mathbb{I}}
\newcommand{\R}{\mathbb{R}}
\newcommand{\Q}{\mathbb{Q}}
\newcommand{\C}{\mathbb{C}}
\newcommand{\Z}{\mathbb{Z}}
\newcommand{\N}{\mathbb{N}}
\newcommand{\xopt}{x^{*}}
\newcommand{\yopt}{y^{*}}
\DeclareMathOperator{\im}{Im}

\newcommand{\nombremateria}
{Investigaci\'on Operativa}
\newcommand{\nombrecuatrimestre}
    {Segundo cuatrimestre de 2025}
\newcommand{\numeroparcial}
    { Segundo Parcial }
\newcommand{\fecha}
    {3 de Octubre de 2025}
\newcommand{\numerotema}
    {\bf 1}

%\oddsidemargin -1cm \evensidemargin -1cm \textwidth 17.5cm
%topmargin 1.2cm \textheight 25cm
%\setlength{\textheight}{25cm}

\def \ds{\displaystyle}
\theoremstyle{definition}
\newtheorem{quest}{Ejercicio}
\thispagestyle{empty}

\begin{document}

\hfill \hfill
\begin{tabular}[c]{|c|c|c|c|}
\hline 
1 & 2 & 3 & Calificaci\'on \\ \hline \hline
\hspace{1.2 cm} & \hspace{1.2 cm} & \hspace{1.2 cm} &\hspace{1.5 cm} \\
\hspace{1.2 cm} & \hspace{1.2 cm} & \hspace{1.2 cm} &\hspace{1.5 cm} \\ \hline

\end{tabular}

\vspace{-1cm}

\noindent Nombre y apellido:

%\vspace{0.2cm}

\noindent Número de libreta:

\noindent \hrulefill 

%\smallskip
\renewcommand{\baselinestretch}{1.1}

\begin{center}
	\large \textbf{\nombremateria} 
	
	\numeroparcial -- \fecha
\end{center}

%\vspace{.6cm}
\renewcommand{\theenumi}{\textbf{\arabic{enumi}}}

\begin{center}
    \textbf{Entregar todos los ejercicios en hojas separadas.}
\end{center}
%%% EMPIEZA EL EJERCICIO 1 %%%

\begin{quest} 
\verb|[5.5 pts.]| Dado el siguiente problema de Programación Lineal:
    \[
        \begin{array}{rrcrcrcr}
            \max & 3x_1 & + & 2x_2 & + & 2x_3 & & \\
            \text{s.a:} & x_1 & + & 2x_2 & + & x_3 & \leq & -4\\
            & & - & x_2 & - & x_3 & \leq & -2\\
            & x_1 & + & x_2 &  & & \leq & -1\\
            & \multicolumn{5}{r}{x_1,x_3} & \leq & 0 \\
            & \multicolumn{5}{r}{x_2} &  & \text{libre} \\

        \end{array}
    \]
    \begin{enumerate}[label=\alph*)]
        \item \verb|[1 pt.]| Plantear el problema dual.
        \item \verb|[3 pts.]| Resolver el problema dual con el Método de Dos Fases.
        \item \verb|[1.5 pts.]| A partir de la solución óptima obtenida en el ítem anterior, hallar una solución
        óptima del problema
        primal.
    \end{enumerate}

\end{quest}

%%% TERMINA EL EJERCICIO 1 %%%

%%% EMPIEZA EL EJERCICIO 2 %%%

\begin{quest}
    \verb|[1 pt.]| El siguiente diccionario corresponde a una solución básica factible de un problema de
    Programación Lineal con objetivo de \textbf{maximizar}:
    \[
    \begin{array}{rrrrrrrrr}
        x_4&=&a&-&2x_1&+&x_2&-&\frac{1}{2}x_3\\
        x_5&=&1&-&bx_1&+&x_2&+&cx_3\\
        x_6&=&2&&&-&x_2&-&3x_3\\ \hline
        z&=&&&dx_1&-&2x_2 &+&ex_3
    \end{array}
    \]
    con $a,b,c,d,e\in \mathbb{R}$ y $a\geq 0$. Utilizando el criterio usual de elección de qué variable entra a la
    base, decidir si cada una de las siguientes afirmaciones son verdaderas o falsas y
    \underline{justificar la respuesta}.
    \begin{enumerate}[label=\alph*)]
        \item \verb|[0.25 pts.]| Para que $x_1$ entre a la base y salga $x_5$ basta pedir:
        \[d>e>0 \;\wedge\; b>0 \;\wedge\; \frac{1}{b} < \frac{a}{2} \]
        \item \verb|[0.25 pts.]| Si $a=0$, la próxima iteración dará lugar a una solución degenerada.
        \item \verb|[0.25 pts.]| Si $a\neq 0$, $d=0$ y $e\leq0$, el problema tiene infinitas soluciones óptimas.
        \item \verb|[0.25 pts.]| Si $e>d>0$ y $c > 0$, se puede afirmar que el problema es no acotado.
    \end{enumerate}

\end{quest}
%%% TERMINA EL EJERCICIO 2 %%%

%%% EMPIEZA EL EJERCICIO 3 %%%

\begin{quest} 
\verb|[3.5 pts.]| Utilizar el algoritmo de Branch \& Bound para resolver el siguiente problema de Programación Lineal
Entera, detallando cada división en subproblemas mediante el esquema de árbol visto en clase.
\[
\begin{array}{rrcl}
\max & \multicolumn{3}{l}{x_1+3x_2} \\
s.a: & 2x_1+x_2 & \leq & 8 \\
 & 4x_1+9x_2 & \leq & 24 \\
 & 4x_1-5x_2 & \geq & -4 \\
% & -3x_1+3x_2 & \geq & -11 \\
 & x_1, x_2 & \geq & 0 \\
 & x_1, x_2 & \in & \mathbb{Z}
\end{array}
\]
\end{quest}

%%% TERMINA EL EJERCICIO 3 %%%

%%% TERMINA EL PARCIAL %%%

\newpage

   \begin{tabular}{|l|l|} \hline
        \multicolumn{1}{|c|}{Primal} & \multicolumn{1}{c|}{Dual}  \\ \hline
        Obj: Minimizar & Obj: Maximizar \\ \hline
        % Obj: Minimizar & Obj: Maximizar \\ \hline
        $\imath$-ésima restricción $\leq$ & $\imath$-ésima variable $\leq 0$ \\
        $\imath$-ésima restricción $\geq$ & $\imath$-ésima variable $\geq 0$ \\
        $\imath$-ésima restricción $=$ & $\imath$-ésima variable libre \\
        $\jmath$-ésima variable $\geq 0$ & $\jmath$-ésima restricción $\leq$ \\
        $\jmath$-ésima variable $\leq 0$ & $\jmath$-ésima restricción $\geq$ \\
        $\jmath$-ésima variable libre & $\jmath$-ésima restricción $=$ \\ \hline
    \end{tabular}
    \begin{tabular}{|l|l|} \hline
        \multicolumn{1}{|c|}{Primal} & \multicolumn{1}{c|}{Dual}  \\ \hline
        Obj: Maximizar & Obj: Minimizar \\ \hline
        % Obj: Minimizar & Obj: Maximizar \\ \hline
        $\imath$-ésima restricción $\leq$ & $\imath$-ésima variable $\geq 0$ \\
        $\imath$-ésima restricción $\geq$ & $\imath$-ésima variable $\leq 0$ \\
        $\imath$-ésima restricción $=$ & $\imath$-ésima variable libre \\
        $\jmath$-ésima variable $\geq 0$ & $\jmath$-ésima restricción $\geq$ \\
        $\jmath$-ésima variable $\leq 0$ & $\jmath$-ésima restricción $\leq$ \\
        $\jmath$-ésima variable libre & $\jmath$-ésima restricción $=$ \\ \hline
    \end{tabular}
    \vspace{1cm}

    \textbf{Teorema débil de dualidad:} Para un problema con objetivo de maximizar, sean $(x_1,\cdots,x_n)$ solución
factible del primal e $(y_1,\cdots, y_m)$ solución factible del dual, entonces:
    \[ \displaystyle \sum_{j=1}^n c_jx_j \leq \sum_{i=1}^m b_iy_i\]
    (si el objetivo del primal es minimizar, se invierte la desigualdad).

    \textbf{Teorema fundamental de dualidad}: Si el problema primal tiene solución óptima $(x_1^{\ast},\cdots,x_n^{\ast})$, entonces el dual
    tiene solución óptima $(y_1^{\ast},\cdots,y_m^{\ast})$ tal que:
    \[ \displaystyle \sum_{j=1}^n c_jx_j^{\ast} = \sum_{i=1}^m b_iy_i^{\ast}\]

    \textbf{Teorema de Holgura Complementaria:} Sean $x^{*}=(x_1^{\ast},\cdots,x_n^{\ast})$ solución del primal e
        $y^{*}=(y_1^{\ast},\cdots,y_m^{\ast})$ solución del dual, las siguientes son condiciones necesarias y
suficientes para la optimalidad simultánea de $\xopt$ e $\yopt$:
    \[
    \begin{array}{ccccc}
        \displaystyle \sum_{i=1}^m a_{ij}\yopt_i=c_j & \text{o} & \xopt_j = 0 & \; \text{(o ambos)} \quad & \forall j=1,\cdots,n  \\
        \displaystyle \sum_{j=1}^m a_{ij}\xopt_j=b_i & \text{o} & \yopt_i = 0 & \; \text{(o ambos)} \quad & \forall i=1,\cdots,m  \\
    \end{array} \]
\end{document}

